\scanpage{10}
or again
\[
\exists\ q,r\in\hat{\mathbb Z}^{*}\quad\text{such that}\quad\boxed{gl_1^q\sigma_\infty(g)=l_0^r}
\]
\margin{it is better to write the relation in the form
\[
\sigma_\infty(g^{-1})=l_0^\beta gl_1^\gamma,
\]
we find, on passing to the abelian group $\beta=\gamma$ (i.e. $q=-r$) thus
\[
\sigma_\infty(g^{-1})=l_0^\beta gl_1^\beta
\]}
(Condition which depends only on $g$, not on $p$;

If $g$ is replaced by $g'=l_0^\alpha gl_1^{-\alpha}$, we shall have $\sigma_\infty(g')=l_1^\alpha\sigma_\infty(g)l_0^{-\alpha}$, and the relation
\[
g'l_1^{q'}\sigma_\infty(g')=l_0^{r'}
\]
is written
\[
l_0^\alpha gl_1^{-\alpha}l_1^{q'}l_1^\alpha\sigma_\infty(g)l_0^{-\alpha}=l_0^{r'}
\]
i.e. simplifying
\[
gl_1^{q'}\sigma_\infty(g)=l_0^{r'}
\]
thus $q=q'$, $r=r'$ i.e. $(q,r)$ depend only on the outer automorphism defined by $u$, and not on the \deleted{\unsure{choice of u|choice of this}} of its representative in the normalizer of $L_0$.

In the case of $\tau$, where $g=l_1$, the relation becomes
\[
l_1l_1^ql_0\unsure{=|=:}l_0^r\quad\text{thus}\quad q=-1,\ r=1.
\]
We know that, if we work with discrete groups, in addition to the identity outer automorphism, there is only one other solution of our equations (requiring \inserted{thus $p=\pm1$ \unsure{and|and)}} that $l_0^p$ and $\operatorname{int}(g)l_1^p$ generate $\pi$, in addition to the fundamental relations
\mspage{9}
